\chapter{Conclusion}
\label{ch:conclusion}

The objective of this project was to determine whether the use of a mobile application could serve as an additional tool to help kids with speech impediments perform their speech therapy exercises outside of a clinical environment. The motivation for this was based on the established gap between the guidance and monitoring received during clinical sessions and the inconsistency of home practice, which often has a loose structure and lacks direction.

The project delivered a prototype consisting of two main components:

\begin{itemize}
    \item A Unity-based mobile application that manages speech exercises in a child-friendly, interactive format recommended by an SLP.
    \item A Django-based web dashboard that enables therapists to define exercises, assign them to different patients, and analyze session completion data.
\end{itemize}

Evaluation of the prototype involved conducting an informal presentation to a speech-language pathologist, who expressed approval of the chosen exercise structure and stated that they would be glad to have such a tool in practice. Specifically, they appreciated the convenience of a mobile system as a replacement for paper exercises, which patients tend to lose track of. The modular exercise system, capable of producing custom exercises, proved to be a feasible idea, relevant to the real-world issues that SLPs face in their day-to-day practice.

On the other hand, the prototype is currently a proof of concept and, as such, has many limitations that would need to be addressed before it could become production-ready. Other experimental features, such as facial tracking, show great promise but have not yet been clinically validated. This limits the amount of automated evaluation that could take place.

Despite this, the prototype has demonstrated the project's feasibility in real-world use cases, with practicing SLPs endorsing the idea that, with continued development, this project has the potential to make a real difference in how children with speech impediments engage in home exercise.
